\chapter{Chlamydia Testing}

\section*{What Is This Test?}
Chlamydia testing detects \textit{Chlamydia trachomatis}, an obligate intracellular bacterium that can cause urethritis, cervicitis, pelvic inflammatory disease, epididymitis, proctitis, neonatal conjunctivitis, and neonatal pneumonia. Infection is frequently asymptomatic. Nucleic acid amplification testing (NAAT) is the preferred diagnostic method, with the specimen site selected from the patient's anatomy and sexual exposure.

\section*{Alternative Names}
Chlamydia NAAT; CT PCR; \textit{C. trachomatis} test; chlamydia culture; CT/NG NAAT.

\section*{Principle}
NAAT detects \textit{C. trachomatis} DNA or RNA and is generally more sensitive than culture. Vaginal swabs are preferred for screening people with a vagina; first-catch urine is commonly used for people with a penis. Rectal and pharyngeal testing should be based on reported exposure and current guidance. Assay-cleared specimen types vary by platform and jurisdiction. Culture is specialized and may be used when maximum specificity, antimicrobial research, or a medicolegal indication is required \cite{workowski2021sexually}.

\section*{History}
Culture and antigen tests preceded modern molecular assays. NAATs became the principal diagnostic approach because they permit sensitive testing of noninvasive specimens and extragenital sites when the assay is validated or cleared for that specimen.

\section*{Why This Test Is Ordered}
\begin{itemize}
  \item Screening sexually active women under 25 and older women with increased risk, according to local recommendations.
  \item Screening during pregnancy according to prenatal guidance and risk; pregnant patients with chlamydia require a test of cure.
  \item Evaluation of urethritis, cervicitis, PID, epididymitis, proctitis, or a sexually exposed partner.
  \item Site-specific testing after receptive anal or oral exposure when clinically indicated.
  \item Evaluation of a newborn for suspected chlamydial conjunctivitis or pneumonia using age-appropriate testing.
\end{itemize}

\section*{Primary vs Secondary Test}
NAAT is the primary diagnostic test for chlamydia at validated specimen sites. Serology is not recommended for routine diagnosis of acute urogenital infection. Culture is a secondary or specialized test.

\section*{How This Test Is Performed}
Collect a self-collected or clinician-collected vaginal swab, endocervical swab, first-catch urine, or an extragenital swab according to the patient's anatomy, exposure, assay instructions, and local validation. First-catch urine generally means the first 10--20 mL of the stream. Place the specimen in the manufacturer's transport medium and label the anatomic site.

\section*{What Preparation Is Needed}
Fasting is not required. Follow the collection instructions for the particular assay. Avoid urinating for approximately 1 hour before first-catch urine when the laboratory recommends this. Recent antibiotic exposure should be documented but is not a general reason to defer diagnostic NAAT. A test of cure, when indicated, should not be performed before 4 weeks after treatment because residual nucleic acid may cause a false-positive result.

\section*{Contraindications and Precautions}
There are no absolute contraindications to urine or swab collection. A negative result applies only to the site tested and can be falsely negative if the specimen is inadequate, collected too early, or obtained after treatment. Routine pharyngeal screening is not generally recommended, but a positive pharyngeal result found during testing should be managed according to current guidance. Test of cure is recommended in pregnancy and may be indicated for persistent symptoms, suspected nonadherence, or selected LGV cases; it is not routine after uncomplicated treatment in nonpregnant patients. Retesting for reinfection is recommended approximately 3 months after treatment.

\section*{Risks}
Urine collection is noninvasive. Swab collection may cause brief discomfort or spotting. The major risks are from untreated infection, including PID, infertility, ectopic pregnancy, epididymitis, and neonatal infection.

\section*{What Happens After the Test}
Results may be available within hours or may take longer depending on the laboratory. A positive result requires treatment and partner management according to current local guidance. Patients should be tested for other relevant STIs, and partners from the preceding 60 days should be evaluated and treated. Patients should avoid sexual activity until the recommended abstinence interval is complete and partners have been treated.

\section*{Understanding Results}
\subsection*{Below Normal}
\begin{itemize}
  \item \textbf{Not detected:} No chlamydial nucleic acid was detected at the tested site. This does not exclude infection at an untested site or a poorly collected specimen.
  \item \textbf{Culture negative:} No organism was recovered; sensitivity is lower than NAAT.
\end{itemize}

\subsection*{Above Normal}
\begin{itemize}
  \item \textbf{Detected:} Chlamydia was detected at the tested site and should be treated according to current guidance.
  \item \textbf{Rectal infection with suspected LGV:} Severe proctitis or proctocolitis with a positive rectal NAAT requires specialist/public-health guidance and, where available, genotyping.
\end{itemize}

\section*{Normal Results}
\textit{C. trachomatis} not detected at each site tested.

\section*{Follow-up Tests}
\begin{itemize}
  \item Retest for reinfection approximately 3 months after treatment.
  \item Test of cure at 4 weeks in pregnancy or when otherwise indicated.
  \item Gonorrhea, HIV, syphilis, and other STI testing according to exposure and risk.
  \item Specialist evaluation and LGV testing for compatible symptomatic rectal infection.
\end{itemize}

\section*{References}
\bibliographystyle{unsrtnat}
\bibliography{references}
\clearpage
\section*{Regulatory Status and Authorization Date}
Regulatory status applies to the specific assay, instrument, manufacturer, and intended use rather than to the test name alone. No single approval date can be assigned to this test category. For a commercial assay, verify the current product in the relevant regulator's database: the U.S. Food and Drug Administration (FDA) clearance, approval, or authorization; India's Central Drugs Standard Control Organisation (CDSCO) licence or permission; the European Union In Vitro Diagnostic Regulation (EU IVDR) CE certificate issued through the applicable conformity-assessment route; or the United Kingdom Medicines and Healthcare products Regulatory Agency (MHRA) registration and UKCA/CE status. Laboratory-developed tests may not have an individual FDA, CDSCO, EU IVDR, or MHRA approval date.
\clearpage
